\documentclass[11pt,a4paper]{article}

\usepackage{../styles/neet_style}

\begin{document}

% --- Cover / Header Banner ---
\begin{center}
  {\huge\bfseries\color{primarydark} NEET UG 2027: Master Examination Syllabus}\\[6pt]
  {\Large\bfseries\color{accentblue} Comprehensive 50-Unit Curriculum Specification \& Exclusion Blueprint}\\[8pt]
  {\normalsize\textbf{Prepared, Compiled \& Authored by Satya Prakash} \quad\textbar\quad \textbf{Target Exam:} National Eligibility cum Entrance Test (NEET UG 2027)}\\[4pt]
  {\footnotesize\color{gray} Strictly Calibrated to National Medical Commission (NMC) Guidelines \& Rationalized NCERT Textbooks}
\end{center}
\vspace{-2pt}
\hrule height 1.2pt
\vspace{8pt}

% --- Executive Summary / Exam Structure ---
\section*{Official Examination Structure \& Marking Architecture}
\addcontentsline{toc}{section}{Official Examination Structure \& Marking Architecture}

The National Eligibility cum Entrance Test (NEET UG 2027) is administered by the National Testing Agency (NTA) based on the regulatory curriculum prescribed by the Under Graduate Medical Education Board (UGMEB) of the National Medical Commission (NMC).

\begin{table}[htbp]
\centering
\small
\renewcommand{\arraystretch}{1.2}
\begin{tabularx}{\textwidth}{l c c c Y}
\toprule
\textbf{Subject Domain} & \textbf{Compulsory Questions} & \textbf{Marks} & \textbf{Target Pace} & \textbf{Question Typology \& Format} \\
\midrule
\textbf{Physics} & 45 (All Compulsory) & 180 & 45--50 min & Single-Choice MCQs, Numerical Calculations, \\
\textbf{Chemistry} & 45 (All Compulsory) & 180 & 40--45 min & Assertion-Reason, Statement I/II Analysis, \\
\textbf{Botany} & 45 (All Compulsory) & 180 & 35--40 min & Diagrammatic Evaluations, \\
\textbf{Zoology} & 45 (All Compulsory) & 180 & 35--40 min & Multi-Column Matrix Matching \\
\midrule
\textbf{Total Examination} & \textbf{180 Questions} & \textbf{720} & \textbf{180 min (3.0 h)} & \textbf{Zero Optional Questions / No Section B} \\
\bottomrule
\end{tabularx}
\caption{Official Reformed NEET UG Examination Scheme (+4 for correct response, $-1$ for incorrect response, 0 for unattempted. Optional Section B discontinued by NTA).}
\end{table}

\begin{trapbox}[NMC Rationalization Warning: Avoid Obsolete Study Modules]
The NMC aligned the NEET UG examination directly with rationalized NCERT textbooks. Legacy coaching material contains 9 entire chapters that are \textbf{completely eliminated} from the syllabus. Studying these chapters wastes hundreds of hours of high-value preparation time. This document flags all excluded and newly prioritized topics explicitly.
\end{trapbox}

\vspace{5pt}
\tableofcontents
\vspace{10pt}
\newpage

% =========================================================================
% SECTION 1: PHYSICS SYLLABUS
% =========================================================================
\section{Physics: Complete 20-Unit Detailed Breakdown}

The Physics curriculum comprises 20 official units spanning Class 11 and Class 12 mechanics, thermodynamics, electrodynamics, optics, modern physics, and a mandatory experimental skills unit.

\subsection{Unit 1: Physics and Measurement}
\begin{itemize}
  \item \textbf{Scope:} Units of measurement, Systems of units (SI Units, fundamental and derived units).
  \item \textbf{Precision \& Measurement:} Least count, significant figures, rounding off.
  \item \textbf{Error Analysis:} Absolute error, relative error, percentage error, combination and propagation of errors in addition, subtraction, multiplication, division, and exponentiation.
  \item \textbf{Dimensional Analysis:} Dimensions of physical quantities, dimensional formulae and equations, dimensional consistency of equations, conversion of units using dimensional method, derivation of relations, limitations of dimensional analysis.
  \item \textbf{Practical Tools:} Vernier Callipers (least count, zero error, zero correction), Screw Gauge (pitch, least count, zero error).
\end{itemize}

\subsection{Unit 2: Kinematics}
\begin{itemize}
  \item \textbf{Motion in a Straight Line:} Frame of reference, position-time coordinate representation, distance and displacement, average and instantaneous speed and velocity, uniformly accelerated motion, graphical kinematics ($x-t$, $v-t$, $a-t$ graphs), kinematic equations of motion for constant acceleration ($v=u+at$, $s=ut+\frac{1}{2}at^2$, $v^2=u^2+2as$, $s_n=u+\frac{a}{2}(2n-1)$), elementary calculus applications ($\frac{dx}{dt}, \frac{dv}{dt}, v\frac{dv}{dx}$).
  \item \textbf{Motion in a Plane:} Scalar and vector quantities, position and displacement vectors, equality of vectors, multiplication of vectors by real numbers, addition and subtraction of vectors (triangle and parallelogram law), relative velocity in 2D (river-boat problems, rain-man problems).
  \item \textbf{Vector Algebra:} Unit vector, resolution of a vector in a plane, rectangular components, scalar (dot) product and properties, vector (cross) product and properties.
  \item \textbf{Projectile Motion:} Motion under gravity in a plane, equation of trajectory, time of flight, maximum height, horizontal range, complementary angles ($45^\circ \pm \theta$).
  \item \textbf{Uniform Circular Motion:} Angular displacement, angular velocity, angular acceleration, centripetal acceleration ($a_c = \frac{v^2}{r} = \omega^2 r$), tangential acceleration.
\end{itemize}

\subsection{Unit 3: Laws of Motion}
\begin{itemize}
  \item \textbf{Foundations:} Intuitive concept of force, inertia, Newton's first law of motion.
  \item \textbf{Dynamics:} Momentum, Newton's second law of motion ($\mathbf{F}=\frac{d\mathbf{p}}{dt}$), impulse of a force ($J = \int F dt = \Delta p$), Newton's third law of motion (action-reaction pairs).
  \item \textbf{Conservation Laws:} Law of conservation of linear momentum and its applications (recoil of gun, rocket propulsion basics).
  \item \textbf{Equilibrium of Concurrent Forces:} Lami's theorem, free-body diagrams (FBD), tension in connected strings and pulley systems.
  \item \textbf{Friction:} Origin of friction, static friction, limiting friction, kinetic friction, laws of friction, coefficient of friction ($\mu_s \ge \mu_k$), angle of friction, angle of repose, motion of bodies on rough horizontal and inclined planes, lubrication.
  \item \textbf{Dynamics of Circular Motion:} Centripetal force, motion of a car on a level circular road ($v_{\max} = \sqrt{\mu g r}$), motion of a car on a banked road (with and without friction, $v_{\text{opt}} = \sqrt{rg\tan\theta}$), conical pendulum.
\end{itemize}

\subsection{Unit 4: Work, Energy and Power}
\begin{itemize}
  \item \textbf{Work:} Work done by a constant force and a variable force ($\int \mathbf{F}\cdot d\mathbf{r}$), graphical interpretation of work (area under $F-x$ curve).
  \item \textbf{Energy:} Kinetic energy, Work-Energy theorem for constant and variable forces ($\Delta K = W_{\text{net}}$).
  \item \textbf{Potential Energy:} Concept of potential energy, conservative and non-conservative forces, potential energy of a spring ($U = \frac{1}{2}kx^2$), conservative force-potential relation ($F = -\frac{dU}{dx}$), points of stable, unstable, and neutral equilibrium.
  \item \textbf{Conservation Principles:} Conservation of mechanical energy, vertical circular motion (critical velocities at bottom $\sqrt{5gl}$, top $\sqrt{gl}$, horizontal $\sqrt{3gl}$, string tension analysis).
  \item \textbf{Power:} Average power, instantaneous power ($P = \mathbf{F}\cdot\mathbf{v}$).
  \item \textbf{Collisions:} Elastic and inelastic collisions in one and two dimensions, coefficient of restitution ($e$), head-on elastic collision formulas, loss of kinetic energy in perfectly inelastic collision.
\end{itemize}

\subsection{Unit 5: Rotational Motion}
\begin{itemize}
  \item \textbf{Centre of Mass:} Centre of mass of a two-particle system, momentum conservation and centre of mass motion, centre of mass of a rigid body, continuous mass distribution (rod, ring, disc, sphere).
  \item \textbf{Rotational Kinematics \& Dynamics:} Basic concepts of rotational motion, torque ($\boldsymbol{\tau} = \mathbf{r}\times\mathbf{F}$), angular momentum ($\mathbf{L} = \mathbf{r}\times\mathbf{p} = I\boldsymbol{\omega}$), relation between torque and angular momentum ($\boldsymbol{\tau} = \frac{d\mathbf{L}}{dt}$), conservation of angular momentum and its applications.
  \item \textbf{Moment of Inertia:} Definition ($I = \sum m_i r_i^2$), radius of gyration ($k$), parallel axis theorem ($I = I_{\text{cm}} + Md^2$), perpendicular axis theorem ($I_z = I_x + I_y$).
  \item \textbf{Standard Values of $I$:} Thin ring, uniform circular disc, thin rod, solid cylinder, hollow cylinder, solid sphere, hollow spherical shell.
  \item \textbf{Rolling Motion:} Pure rolling on horizontal and inclined surfaces, kinetic energy of rolling body ($K = \frac{1}{2}Mv^2(1 + \frac{k^2}{R^2})$), acceleration down an inclined plane ($a = \frac{g\sin\theta}{1 + k^2/R^2}$).
\end{itemize}

\subsection{Unit 6: Gravitation}
\begin{itemize}
  \item \textbf{Planetary Laws:} Kepler's laws of planetary motion (law of orbits, law of areas, law of periods $T^2 \propto r^3$).
  \item \textbf{Universal Gravitation:} Newton's universal law of gravitation, gravitational constant $G$.
  \item \textbf{Acceleration Due to Gravity ($g$):} Acceleration at Earth surface, variation with altitude ($g_h = g(1-\frac{2h}{R})$ for $h \ll R$), variation with depth ($g_d = g(1-\frac{d}{R})$), variation with latitude due to rotation ($g' = g - \omega^2 R \cos^2\lambda$).
  \item \textbf{Gravitational Potential \& Energy:} Gravitational potential energy ($U = -\frac{GMm}{r}$), gravitational potential ($V = -\frac{GM}{r}$), relation between field and potential ($E = -\frac{dV}{dr}$).
  \item \textbf{Satellites \& Orbital Motion:} Escape velocity from earth surface ($v_e = \sqrt{\frac{2GM}{R}} = \sqrt{2gR} \approx 11.2\text{ km/s}$), orbital velocity ($v_o = \sqrt{\frac{GM}{r}}$), relation $v_e = \sqrt{2}v_o$, time period of satellite, total energy of satellite ($E = -K = \frac{U}{2} = -\frac{GMm}{2r}$), geostationary and polar satellites.
\end{itemize}

\subsection{Unit 7: Properties of Solids and Liquids}
\begin{itemize}
  \item \textbf{Elasticity:} Elastic behaviour, stress and strain, Hooke's law, stress-strain diagram, elastic moduli (Young's modulus $Y$, Bulk modulus $B$, Shear modulus $\eta$), compressibility ($K = 1/B$), Poisson's ratio ($\sigma$), elastic potential energy density ($u = \frac{1}{2}\times\text{stress}\times\text{strain}$).
  \item \textbf{Fluid Statics:} Pressure in fluid column, Pascal's law and its applications (hydraulic lift, hydraulic brakes), effect of gravity on fluid pressure, buoyant force and Archimedes' principle.
  \item \textbf{Fluid Dynamics:} Viscosity, Newton's law of viscosity, coefficient of viscosity, Stokes' law, terminal velocity ($v_t = \frac{2}{9}\frac{r^2(\rho-\sigma)g}{\eta}$), streamline and turbulent flow, critical velocity, Reynolds number, equation of continuity ($A_1 v_1 = A_2 v_2$), Bernoulli's theorem and its applications (Venturimeter, Torricelli's law of efflux $v=\sqrt{2gh}$, aerodynamic lift).
  \item \textbf{Surface Tension:} Surface energy and surface tension, angle of contact, excess pressure inside a liquid drop ($\Delta P = \frac{2T}{R}$), soap bubble ($\Delta P = \frac{4T}{R}$), capillary rise ($h = \frac{2T\cos\theta}{r\rho g}$), detergent action.
  \item \textbf{Thermal Properties:} Heat, temperature, thermal expansion (linear $\alpha$, areal $\beta$, volumetric $\gamma$), specific heat capacity, molar heat capacity, calorimetry, principle of mixtures, change of state, latent heat of fusion and vaporization.
  \item \textbf{Heat Transfer:} Modes of heat transfer: Conduction (thermal conductivity $k$, temperature gradient), Convection, Radiation; Newton's law of cooling, Stefan-Boltzmann law ($E = \sigma e A T^4$), Wien's displacement law ($\lambda_m T = b$).
\end{itemize}

\subsection{Unit 8: Thermodynamics}
\begin{itemize}
  \item \textbf{Foundational Laws:} Thermal equilibrium, Zeroth law of thermodynamics and concept of temperature.
  \item \textbf{First Law of Thermodynamics:} Heat, work, and internal energy; First law statement ($\Delta Q = \Delta U + \Delta W$), sign conventions.
  \item \textbf{Thermodynamic Processes:} Quasi-static process, Isothermal process ($PV=\text{const}$, $W = nRT\ln\frac{V_2}{V_1}$, $\Delta U=0$), Adiabatic process ($PV^\gamma=\text{const}$, $W = \frac{P_1V_1-P_2V_2}{\gamma-1}$, $\Delta Q=0$), Isobaric process ($W = P\Delta V$), Isochoric process ($W=0, \Delta Q = \Delta U = nC_v\Delta T$), Cyclic processes (work done = area of $P-V$ loop).
  \item \textbf{Second Law of Thermodynamics:} Reversible and irreversible processes, Clausius and Kelvin-Planck statements.
\end{itemize}

\subsection{Unit 9: Kinetic Theory of Gases}
\begin{itemize}
  \item \textbf{Ideal Gas Equation:} Equation of state of a perfect gas ($PV = nRT = Nk_B T$).
  \item \textbf{Postulates \& Pressure:} Kinetic postulates, derivation of pressure exerted by an ideal gas ($P = \frac{1}{3}\rho v_{\text{rms}}^2$).
  \item \textbf{Kinetic Interpretation of Temperature:} Average kinetic energy per molecule ($\frac{3}{2}k_B T$), RMS speed ($v_{\text{rms}} = \sqrt{\frac{3RT}{M}}$), average speed ($v_{\text{avg}} = \sqrt{\frac{8RT}{\pi M}}$), most probable speed ($v_{\text{mp}} = \sqrt{\frac{2RT}{M}}$).
  \item \textbf{Degrees of Freedom:} Law of equipartition of energy, degrees of freedom for monoatomic ($f=3$), diatomic ($f=5$ at room temp), and polyatomic gases; calculation of $C_v, C_p$, and $\gamma$.
  \item \textbf{Transport Concepts:} Mean free path ($\lambda = \frac{1}{\sqrt{2}\pi n d^2}$), Avogadro's number.
\end{itemize}

\subsection{Unit 10: Oscillations and Waves}
\begin{itemize}
  \item \textbf{Periodic Motion \& SHM:} Periodic and oscillatory motion, period, frequency, displacement function, Simple Harmonic Motion (SHM) differential equation ($\frac{d^2 x}{dt^2} + \omega^2 x = 0$), phase and initial phase constant.
  \item \textbf{Kinematics and Energy in SHM:} Displacement $x(t) = A\sin(\omega t + \phi)$, velocity $v(t) = \omega\sqrt{A^2-x^2}$, acceleration $a(t) = -\omega^2 x$; Kinetic energy $K = \frac{1}{2}m\omega^2(A^2-x^2)$, potential energy $U = \frac{1}{2}m\omega^2 x^2$, total mechanical energy $E = \frac{1}{2}m\omega^2 A^2$ (constant).
  \item \textbf{Oscillating Systems:} Spring-block system (horizontal and vertical, $T = 2\pi\sqrt{\frac{m}{k}}$, series combination $\frac{1}{k_s} = \frac{1}{k_1} + \frac{1}{k_2}$, parallel combination $k_p = k_1 + k_2$); Simple pendulum ($T = 2\pi\sqrt{\frac{L}{g}}$).
  \item \textbf{Damping \& Resonance:} Free, damped, and forced oscillations, resonance.
  \item \textbf{Wave Motion:} Longitudinal and transverse waves, progressive wave equation ($y = A\sin(kx - \omega t + \phi)$), wave speed $v = \frac{\omega}{k} = \nu\lambda$; speed of transverse waves on stretched string ($v = \sqrt{\frac{T}{\mu}}$), speed of longitudinal wave in fluids ($v = \sqrt{\frac{B}{\rho}}$), Newton-Laplace formula for sound in gas ($v = \sqrt{\frac{\gamma P}{\rho}}$).
  \item \textbf{Superposition \& Standing Waves:} Principle of superposition, reflection of waves at fixed and free ends; standing waves on strings (nodes, antinodes, frequencies $\nu_n = \frac{nv}{2L}$); standing waves in organ pipes (closed organ pipe $\nu_n = (2n-1)\frac{v}{4L}$, open organ pipe $\nu_n = \frac{nv}{2L}$); Beats ($\nu_{\text{beat}} = |\nu_1 - \nu_2|$).
\end{itemize}

\subsection{Unit 11: Electrostatics}
\begin{itemize}
  \item \textbf{Electric Charges \& Fields:} Electric charges, conservation and quantization of charge ($q = ne$), Coulomb's law in scalar and vector form, superposition principle.
  \item \textbf{Electric Field:} Electric field due to a point charge, field lines and properties, electric dipole, electric dipole moment ($\mathbf{p} = q\cdot 2\mathbf{a}$), electric field due to dipole on axial line ($E_{\text{axial}} \approx \frac{2kp}{r^3}$) and equatorial line ($E_{\text{eq}} \approx \frac{kp}{r^3}$), torque on a dipole in a uniform electric field ($\boldsymbol{\tau} = \mathbf{p}\times\mathbf{E}$), potential energy of dipole ($U = -\mathbf{p}\cdot\mathbf{E}$).
  \item \textbf{Gauss's Law:} Electric flux ($\Phi = \oint \mathbf{E}\cdot d\mathbf{A}$), Gauss's theorem statement; applications: field due to infinitely long straight wire ($E = \frac{\lambda}{2\pi\varepsilon_0 r}$), uniformly charged infinite plane sheet ($E = \frac{\sigma}{2\varepsilon_0}$), thin spherical shell (inside $E=0$, outside $E = \frac{kq}{r^2}$).
  \item \textbf{Electric Potential:} Potential difference, electric potential due to point charge, electric dipole, system of charges; equipotential surfaces and properties; electrical potential energy of two-point charge system.
  \item \textbf{Conductors \& Capacitors:} Conductors in electrostatic field, dielectric materials and polarization, capacitance ($C = \frac{Q}{V}$), parallel plate capacitor with and without dielectric ($C = \frac{\kappa \varepsilon_0 A}{d}$), capacitors in series and parallel combinations, energy stored in capacitor ($U = \frac{1}{2}CV^2 = \frac{Q^2}{2C} = \frac{1}{2}QV$), energy density ($u = \frac{1}{2}\varepsilon_0 E^2$).
\end{itemize}

\subsection{Unit 12: Current Electricity}
\begin{itemize}
  \item \textbf{Microscopic Current:} Electric current, drift velocity ($v_d = \frac{eE\tau}{m}$), relaxation time, relation between current and drift velocity ($I = n e A v_d$), mobility ($\mu = \frac{v_d}{E}$).
  \item \textbf{Ohm's Law \& Resistance:} Ohm's law, $V-I$ characteristics, electrical resistance, resistivity ($\rho = \frac{m}{n e^2 \tau}$), conductivity ($\sigma = \frac{1}{\rho}$), temperature dependence of resistance ($R_T = R_0(1+\alpha\Delta T)$).
  \item \textbf{DC Circuits \& Cells:} Internal resistance of a cell, terminal potential difference and EMF ($V = \mathcal{E} - Ir$ during discharge, $V = \mathcal{E} + Ir$ during charging), combination of cells in series and parallel.
  \item \textbf{Circuit Analysis:} Kirchhoff's first law (junction rule) and second law (loop rule).
  \item \textbf{Measurement Instruments:} Wheatstone bridge balance condition ($\frac{P}{Q} = \frac{R}{S}$), Metre bridge for unknown resistance measurement.
\end{itemize}

\subsection{Unit 13: Magnetic Effects of Current and Magnetism}
\begin{itemize}
  \item \textbf{Magnetic Field Laws:} Concept of magnetic field, Oersted's experiment, Biot-Savart law in vector form ($d\mathbf{B} = \frac{\mu_0}{4\pi}\frac{I d\mathbf{l}\times\hat{\mathbf{r}}}{r^2}$), magnetic field on the axis and centre of a circular current loop ($B_{\text{centre}} = \frac{\mu_0 I}{2R}$, $B_{\text{axis}} = \frac{\mu_0 I R^2}{2(R^2+x^2)^{3/2}}$), Ampere's circuital law and application to straight wire and long solenoid ($B = \mu_0 n I$).
  \item \textbf{Lorentz Force:} Force on a moving charge in uniform magnetic field ($\mathbf{F} = q(\mathbf{v}\times\mathbf{B})$), motion in magnetic field (radius $r = \frac{mv}{qB}$, pitch), Lorentz force ($\mathbf{F} = q(\mathbf{E} + \mathbf{v}\times\mathbf{B})$), velocity selector.
  \item \textbf{Conductor in Magnetic Field:} Force on a current-carrying conductor in a uniform magnetic field ($\mathbf{F} = I(\mathbf{L}\times\mathbf{B})$), force between two parallel current-carrying conductors ($\frac{F}{L} = \frac{\mu_0 I_1 I_2}{2\pi d}$), standard definition of ampere.
  \item \textbf{Torque on Loop \& Galvanometer:} Torque on a current loop ($\boldsymbol{\tau} = \mathbf{M}\times\mathbf{B}$, where $\mathbf{M} = NI\mathbf{A}$), Moving Coil Galvanometer (deflection $\theta \propto I$), current sensitivity and voltage sensitivity, conversion of galvanometer to ammeter (shunt $S$) and voltmeter (high series resistance $R$).
  \item \textbf{Magnetism \& Matter:} Current loop as magnetic dipole, bar magnet as equivalent solenoid, magnetic field lines; magnetic dipole in uniform magnetic field (torque and potential energy); magnetic materials: diamagnetic, paramagnetic, and ferromagnetic substances with microscopic origin, Curie's law, susceptibility ($\chi$) and relative permeability ($\mu_r$).
\end{itemize}

\subsection{Unit 14: Electromagnetic Induction and Alternating Currents}
\begin{itemize}
  \item \textbf{Electromagnetic Induction:} Magnetic flux ($\Phi = \mathbf{B}\cdot\mathbf{A}$), Faraday's laws of induction, induced EMF and current, Lenz's law and energy conservation, motional EMF ($\mathcal{E} = Bvl$), Eddy currents.
  \item \textbf{Inductance:} Self-inductance ($L$, $\mathcal{E} = -L\frac{dI}{dt}$), energy stored in an inductor ($U = \frac{1}{2}LI^2$), mutual inductance ($M$, $\mathcal{E}_2 = -M\frac{dI_1}{dt}$).
  \item \textbf{Alternating Current:} Alternating currents, peak, average, and RMS values ($I_{\text{rms}} = \frac{I_0}{\sqrt{2}}$), AC circuits with pure resistor, pure inductor, and pure capacitor; inductive reactance ($X_L = \omega L$), capacitive reactance ($X_C = \frac{1}{\omega C}$).
  \item \textbf{Series LCR Circuit:} Phasor diagram, impedance ($Z = \sqrt{R^2 + (X_L - X_C)^2}$), phase angle ($\tan\phi = \frac{X_L - X_C}{R}$), resonance frequency ($\omega_0 = \frac{1}{\sqrt{LC}}$), quality factor ($Q = \frac{\omega_0 L}{R} = \frac{1}{R}\sqrt{\frac{L}{C}}$), power in AC circuits ($P = V_{\text{rms}}I_{\text{rms}}\cos\phi$), power factor ($\cos\phi = \frac{R}{Z}$), wattless current.
  \item \textbf{AC Devices:} AC generator, Transformer (principle, turns ratio $\frac{V_s}{V_p} = \frac{N_s}{N_p} = \frac{I_p}{I_s}$, energy losses in transformers).
\end{itemize}

\subsection{Unit 15: Electromagnetic Waves}
\begin{itemize}
  \item \textbf{Displacement Current:} Ampere-Maxwell law, need for displacement current ($I_d = \varepsilon_0 \frac{d\Phi_E}{dt}$).
  \item \textbf{Characteristics of EM Waves:} Transverse nature, electric and magnetic vectors perpendicular to each other and direction of propagation ($\mathbf{S} = \frac{1}{\mu_0}\mathbf{E}\times\mathbf{B}$), speed of EM waves ($c = \frac{1}{\sqrt{\mu_0\varepsilon_0}} = \frac{E_0}{B_0}$), momentum and radiation pressure.
  \item \textbf{Electromagnetic Spectrum:} Radio waves, microwaves, infrared, visible, ultraviolet, X-rays, gamma rays; frequency/wavelength ranges, production, detection, and practical applications.
\end{itemize}

\subsection{Unit 16: Optics (Ray and Wave)}
\begin{itemize}
  \item \textbf{Ray Optics - Reflection \& Refraction:} Reflection at spherical mirrors, mirror formula ($\frac{1}{v} + \frac{1}{u} = \frac{1}{f}$), magnification; refraction of light, Snell's law, absolute and relative refractive index, lateral shift, apparent depth; total internal reflection (TIR), critical angle ($\sin i_c = \frac{1}{\mu}$), optical fibres, mirage.
  \item \textbf{Lenses \& Prisms:} Refraction at spherical surfaces ($\frac{\mu_2}{v} - \frac{\mu_1}{u} = \frac{\mu_2 - \mu_1}{R}$), Lens Maker's formula ($\frac{1}{f} = (\mu - 1)(\frac{1}{R_1} - \frac{1}{R_2})$), thin lens formula ($\frac{1}{v} - \frac{1}{u} = \frac{1}{f}$), power of lens ($P = \frac{1}{f}$), combination of thin lenses in contact ($P = P_1 + P_2$); refraction through prism, angle of deviation, minimum deviation formula ($\mu = \frac{\sin(\frac{A + \delta_m}{2})}{\sin(\frac{A}{2})}$), dispersion through prism.
  \item \textbf{Optical Instruments:} Magnifying power of simple microscope, compound microscope ($m = -\frac{v_0}{u_0}(1 + \frac{D}{f_e})$ or $m = -\frac{L}{f_0}\frac{D}{f_e}$), astronomical telescope (refracting and reflecting, $m = -\frac{f_0}{f_e}$, tube length $L = f_0 + f_e$).
  \item \textbf{Wave Optics:} Wave front, Huygens' principle, proof of laws of reflection and refraction using Huygens' principle; interference of light, coherent sources, Young's Double Slit Experiment (YDSE), path difference, condition for constructive ($\Delta x = n\lambda$) and destructive ($\Delta x = (2n-1)\frac{\lambda}{2}$) interference, fringe width ($\beta = \frac{\lambda D}{d}$), shift due to thin transparent sheet; diffraction due to single slit, central maximum width ($\frac{2\lambda D}{a}$), secondary maxima and minima conditions.
\end{itemize}

\subsection{Unit 17: Dual Nature of Matter and Radiation}
\begin{itemize}
  \item \textbf{Photoelectric Effect:} Hertz and Lenard's observations, experimental setup, effect of intensity, frequency, and potential; stopping potential ($eV_0 = K_{\max}$).
  \item \textbf{Einstein's Photoelectric Equation:} Photon hypothesis ($E = h\nu$), work function ($\Phi_0 = h\nu_0$), equation ($K_{\max} = h\nu - \Phi_0 = eV_0$), slopes of $V_0-\nu$ graphs.
  \item \textbf{Matter Waves:} Dual nature of matter, de Broglie relation ($\lambda = \frac{h}{p} = \frac{h}{mv}$), wavelength of charged particles accelerated through potential $V$ ($\lambda = \frac{h}{\sqrt{2mqV}}$, for electron $\lambda = \frac{1.227}{\sqrt{V}}\text{ nm}$).
\end{itemize}

\subsection{Unit 18: Atoms and Nuclei}
\begin{itemize}
  \item \textbf{Atomic Models:} Alpha-particle scattering experiment, Rutherford's planetary model and its limitations; Bohr's postulates for hydrogen atom, radius of $n$-th orbit ($r_n = 0.529 \frac{n^2}{Z}\text{ \AA}$), velocity ($v_n = 2.18\times 10^6 \frac{Z}{n}\text{ m/s}$), energy levels ($E_n = -13.6 \frac{Z^2}{n^2}\text{ eV}$); Hydrogen spectrum (Lyman, Balmer, Paschen, Brackett, Pfund series, Rydberg formula $\frac{1}{\lambda} = R Z^2(\frac{1}{n_1^2} - \frac{1}{n_2^2})$).
  \item \textbf{Nuclear Physics:} Composition and size of nucleus ($R = R_0 A^{1/3}$), nuclear density (independent of mass number $A$), atomic masses, mass defect ($\Delta m$), mass-energy equivalence ($E = \Delta m c^2$), binding energy per nucleon curve and its characteristics; nuclear forces (short range, strong, charge-independent, saturated); nuclear fission and fusion.
\end{itemize}

\subsection{Unit 19: Electronic Devices}
\begin{itemize}
  \item \textbf{Semiconductor Physics:} Energy bands in conductors, semiconductors, and insulators; intrinsic semiconductors ($n_e = n_h = n_i$), extrinsic semiconductors: n-type (donor impurities) and p-type (acceptor impurities), mass action law ($n_e n_h = n_i^2$).
  \item \textbf{Semiconductor Diode:} p-n junction formation, depletion region, barrier potential; forward bias and reverse bias $I-V$ characteristics; diode as a half-wave and full-wave rectifier (efficiency, ripple factor).
  \item \textbf{Special Purpose Diodes:} Light Emitting Diode (LED), Photodiode (reverse bias operation), Solar Cell.
  \item \textbf{Logic Gates:} Digital signals, Boolean operations, truth tables for fundamental gates (OR, AND, NOT) and universal gates (NAND, NOR).
\end{itemize}

\subsection{Unit 20: Experimental Skills (NMC Explicit Practical Unit)}
\begin{itemize}
  \item \textbf{Mechanics Experiments:}
  \begin{enumerate}
    \item Vernier Callipers: Internal/external diameter and depth of spherical/cylindrical objects.
    \item Screw Gauge: Diameter of thin wire and thickness of a thin sheet.
    \item Simple Pendulum: Dissipation of energy by plotting a graph between square of amplitude and time; Determination of acceleration due to gravity $g$.
    \item Metre Scale: Mass of an object by the principle of moments.
    \item Young's Modulus of elasticity of the material of a metallic wire.
    \item Capillarity: Surface tension of water by capillary rise phenomenon.
    \item Viscosity: Terminal velocity of a spherical body in a viscous liquid.
    \item Speed of sound in air at room temperature using a resonance tube.
  \end{enumerate}
  \item \textbf{Thermodynamics Experiments:}
  \begin{enumerate}
    \setcounter{enumi}{8}
    \item Specific heat capacity of a given (i) solid and (ii) liquid by method of mixtures.
  \end{enumerate}
  \item \textbf{Electricity \& Magnetism Experiments:}
  \begin{enumerate}
    \setcounter{enumi}{9}
    \item Metre Bridge: Resistance of a given wire / specific resistance of its material.
    \item Ohm's Law: Verification of Ohm's law and determination of resistivity of a wire.
    \item Galvanometer: Figure of merit and resistance of a galvanometer by half-deflection method; conversion into ammeter/voltmeter.
  \end{enumerate}
  \item \textbf{Optics Experiments:}
  \begin{enumerate}
    \setcounter{enumi}{12}
    \item Focal length of: (i) Convex mirror, (ii) Concave mirror, (iii) Convex lens using $u-v$ method.
    \item Refractive index of a glass slab using a traveling microscope.
    \item Refractive index of a glass prism using angle of minimum deviation graph.
  \end{enumerate}
  \item \textbf{Electronics Experiments:}
  \begin{enumerate}
    \setcounter{enumi}{15}
    \item $I-V$ characteristic curves of a p-n junction diode in forward and reverse bias.
    \item Reverse breakdown characteristic of a Zener diode.
    \item Identification of electronic components (resistor, capacitor, diode, LED, transistor, IC).
  \end{enumerate}
\end{itemize}

\begin{trapbox}[Physics NMC Deletions \& Exclusions]
The following topics are \textbf{permanently excluded} from NEET UG:
\begin{itemize}
  \item Geotextile and obsolete units; Rolling motion complex tensor analysis; Heat engines and refrigerators (Carnot engine); Doppler effect in sound; Van de Graaff generator; Colour coding of resistors; Potentiometer (meter bridge is retained, potentiometer deleted); Cyclotron; Earth's magnetism elements (dip, declination, horizontal component); Resolving power of microscope and telescope; Polarisation of light (Brewster's law, Malus' law); Davisson-Germer experiment; Radioactive decay law ($\alpha, \beta, \gamma$ decays, half life/mean life); Transistors (CE characteristics, amplifier).
\end{itemize}
\end{trapbox}

\newpage

% =========================================================================
% SECTION 2: CHEMISTRY SYLLABUS
% =========================================================================
\section{Chemistry: Complete 20-Unit Detailed Breakdown}

Chemistry is partitioned into Physical Chemistry (8 Units), Inorganic Chemistry (4 Units), and Organic Chemistry (7 Units + 1 Practical Chemistry Unit).

\subsection{Physical Chemistry}

\subsubsection{Unit 1: Some Basic Concepts in Chemistry}
\begin{itemize}
  \item Matter and its nature, Dalton's atomic theory, concept of atom, molecule, element, compound.
  \item Laws of chemical combination (conservation of mass, definite proportions, multiple proportions, Gay-Lussac's law, Avogadro's law).
  \item Atomic and molecular masses, mole concept, molar mass, percentage composition, empirical and molecular formulae.
  \item Chemical equations, stoichiometry and calculations based on stoichiometry, limiting reagent.
  \item Concentration terms: Molarity ($M$), Molality ($m$), Normality ($N$), Mole fraction ($x$), Mass percentage (\% w/w, \% w/v, ppm).
\end{itemize}

\subsubsection{Unit 2: Atomic Structure}
\begin{itemize}
  \item Nature of electromagnetic radiation, photoelectric effect; Spectrum of hydrogen atom, Bohr model of hydrogen atom: postulates, derivation of radius, energy of electron, limitations of Bohr model.
  \item Dual nature of matter, de Broglie's relationship, Heisenberg uncertainty principle ($\Delta x \cdot \Delta p \ge \frac{h}{4\pi}$).
  \item Quantum mechanical model: Concept of atomic orbitals, wave function ($\psi$ and $|\psi|^2$), quantum numbers ($n, l, m_l, m_s$) and their significance.
  \item Shapes of $s, p$, and $d$ orbitals, electron spin and spin quantum number.
  \item Rules for filling electrons in orbitals: Aufbau principle, Pauli's exclusion principle, Hund's rule of maximum multiplicity; electronic configuration of elements, stability of half-filled and completely filled orbitals (exchange energy).
\end{itemize}

\subsubsection{Unit 3: Chemical Bonding and Molecular Structure}
\begin{itemize}
  \item Kossel-Lewis approach, ionic bond, lattice enthalpy, Born-Haber cycle.
  \item Covalent bond: Lewis structures, octet rule and its limitations, formal charge.
  \item Bond parameters: Bond length, bond angle, bond enthalpy, bond order.
  \item Resonance and resonance hybrid, dipole moment and molecular polarity (percentage ionic character).
  \item Valence Shell Electron Pair Repulsion (VSEPR) theory and geometry of molecules ($AX_2$ to $AX_6$, including lone pair distortions).
  \item Valence Bond Theory (VBT): Directional characteristics of covalent bonds, hybridization involving $s, p$, and $d$ orbitals ($sp, sp^2, sp^3, dsp^2, sp^3d, sp^3d^2$), shapes of simple molecules.
  \item Molecular Orbital Theory (MOT): LCAO method, bonding and antibonding orbitals, energy level diagrams for homonuclear diatomic molecules ($H_2$ to $O_2$), bond order, magnetic behaviour ($O_2$ paramagnetism).
  \item Hydrogen bonding: Intermolecular and intramolecular, consequences on physical properties.
\end{itemize}

\subsubsection{Unit 4: Chemical Thermodynamics}
\begin{itemize}
  \item Fundamentals: System and surroundings, open, closed, and isolated systems, extensive and intensive properties, state functions vs path functions.
  \item First Law of Thermodynamics: Work ($w = -P_{\text{ext}}\Delta V$), heat ($q$), internal energy ($U$), $\Delta U = q + w$.
  \item Enthalpy ($H = U + PV$), $\Delta H = \Delta U + \Delta n_g RT$, heat capacities ($C_p, C_v$) and relation $C_p - C_v = R$.
  \item Thermochemistry: Standard states, enthalpy of reaction, enthalpy of formation (Hess's law of constant heat summation), enthalpy of combustion, neutralization, atomization, sublimation, phase transition, ionization, electron gain, bond dissociation enthalpy.
  \item Spontaneity: Criteria for spontaneity, concept of entropy ($S$), second law of thermodynamics, $\Delta S_{\text{sys}} = \frac{q_{\text{rev}}}{T}$, third law of thermodynamics.
  \item Gibbs Free Energy ($G = H - TS$): Standard Gibbs free energy change ($\Delta G^\circ$), criterion for spontaneity ($\Delta G < 0$), relation between $\Delta G^\circ$ and equilibrium constant ($\Delta G^\circ = -RT\ln K$).
\end{itemize}

\subsubsection{Unit 5: Solutions}
\begin{itemize}
  \item Types of solutions, expression of concentration of solutions of liquids in liquids and solids in liquids.
  \item Solubility of gases in liquids: Henry's law and its applications.
  \item Vapour pressure of liquid solutions: Raoult's law for volatile solutes, ideal and non-ideal solutions, positive and negative deviations from Raoult's law, azeotropes (minimum and maximum boiling).
  \item Colligative Properties:
  \begin{enumerate}
    \item Relative lowering of vapour pressure ($\frac{P^\circ - P}{P^\circ} = x_2$).
    \item Elevation of boiling point ($\Delta T_b = K_b m$, where $K_b = \frac{R T_b^2 M}{1000 \Delta H_{\text{vap}}}$).
    \item Depression of freezing point ($\Delta T_f = K_f m$).
    \item Osmotic pressure ($\Pi = CRT$), reverse osmosis and water purification.
  \end{enumerate}
  \item Abnormal Molar Mass: Van 't Hoff factor ($i$), degree of dissociation ($\alpha = \frac{i-1}{n-1}$), degree of association ($\alpha = \frac{1-i}{1 - 1/n}$).
\end{itemize}

\subsubsection{Unit 6: Equilibrium}
\begin{itemize}
  \item \textbf{Chemical Equilibrium:} Dynamic nature of equilibrium, Law of mass action, Equilibrium constants ($K_c$ and $K_p$), relation $K_p = K_c(RT)^{\Delta n_g}$, reaction quotient ($Q$) and direction of reaction, Le Chatelier's principle (effect of concentration, temperature, pressure, addition of inert gas, catalyst).
  \item \textbf{Ionic Equilibrium:} Ionization of electrolytes, strong and weak electrolytes, Arrhenius, Bronsted-Lowry, and Lewis concepts of acids and bases, ionization of water and ionic product ($K_w = 10^{-14}$ at $25^\circ\text{C}$).
  \item \textbf{pH Scale:} Definition, calculation of pH for strong acids, strong bases, and weak acids/bases ($K_a, K_b$, Ostwald's dilution law).
  \item \textbf{Hydrolysis of Salts:} Salts of (strong acid + strong base), (strong acid + weak base), (weak acid + strong base), (weak acid + weak base), hydrolysis constant ($K_h$), degree of hydrolysis ($h$), pH formulas.
  \item \textbf{Buffer Solutions:} Acidic buffer (Henderson-Hasselbalch equation $\text{pH} = \text{p}K_a + \log\frac{[\text{salt}]}{[\text{acid}]}$), basic buffer ($\text{pOH} = \text{p}K_b + \log\frac{[\text{salt}]}{[\text{base}]}$), buffer capacity.
  \item \textbf{Solubility Equilibrium:} Solubility product constant ($K_{sp}$), criteria for precipitation ($Q_{sp} > K_{sp}$), common ion effect on solubility.
\end{itemize}

\subsubsection{Unit 7: Redox Reactions and Electrochemistry}
\begin{itemize}
  \item \textbf{Redox Reactions:} Concept of oxidation and reduction, oxidation number rules, balancing redox reactions (oxidation number method and ion-electron half-reaction method).
  \item \textbf{Electrochemical Cells:} Galvanic cells, Daniel cell, representation of cell, electrode potential, standard electrode potential, SHE (Standard Hydrogen Electrode), Electrochemical series and applications.
  \item \textbf{Nernst Equation:} Derivation and application to electrode potential and cell potential ($E_{\text{cell}} = E^\circ_{\text{cell}} - \frac{0.0591}{n}\log Q$ at $298\text{ K}$), relation between $E^\circ_{\text{cell}}$ and equilibrium constant ($E^\circ_{\text{cell}} = \frac{0.0591}{n}\log K$), $\Delta G^\circ = -n F E^\circ_{\text{cell}}$.
  \item \textbf{Conductance in Solutions:} Resistance, conductance ($G$), resistivity, conductivity ($\kappa = G \cdot \frac{l}{A}$), molar conductivity ($\Lambda_m = \frac{1000\kappa}{M}$), variation of conductivity and molar conductivity with concentration.
  \item \textbf{Kohlrausch's Law:} Independent migration of ions ($\Lambda_m^\circ = \nu_+ \lambda_+^\circ + \nu_- \lambda_-^\circ$), applications (calculation of $\Lambda_m^\circ$ for weak electrolytes, degree of dissociation $\alpha = \frac{\Lambda_m}{\Lambda_m^\circ}$, dissociation constant $K_a$).
  \item \textbf{Electrolysis:} Faraday's laws of electrolysis (first law $m = ZIt$, second law $\frac{m_1}{m_2} = \frac{E_1}{E_2}$), products of electrolysis of molten and aqueous electrolytes.
\end{itemize}

\subsubsection{Unit 8: Chemical Kinetics}
\begin{itemize}
  \item Rate of a chemical reaction: Average and instantaneous rate, factors affecting rate (concentration, temperature, catalyst).
  \item Rate law, rate constant ($k$), order of reaction (zero, first, second, fractional, pseudo-first order), molecularity of reaction, difference between order and molecularity.
  \item Integrated rate equations:
  \begin{itemize}
    \item Zero-order: $[A]_t = [A]_0 - kt$, half-life $t_{1/2} = \frac{[A]_0}{2k}$.
    \item First-order: $k = \frac{2.303}{t}\log\frac{[A]_0}{[A]_t}$, half-life $t_{1/2} = \frac{0.693}{k}$ (independent of initial concentration).
  \end{itemize}
  \item Temperature dependence of rate: Arrhenius equation ($k = A e^{-E_a/RT}$, $\log\frac{k_2}{k_1} = \frac{E_a}{2.303R}[\frac{T_2 - T_1}{T_1 T_2}]$), activation energy ($E_a$), transition state theory, collision theory of chemical reactions (collision frequency $Z$, steric factor $P$).
\end{itemize}

\subsection{Inorganic Chemistry}

\subsubsection{Unit 9: Classification of Elements and Periodicity in Properties}
\begin{itemize}
  \item Modern periodic law and long form of periodic table, electronic configuration and types of elements ($s, p, d, f$ blocks).
  \item Periodic trends in physical properties: Atomic and ionic radii, lanthanoid contraction, isoelectronic species.
  \item Ionization enthalpy: Successive ionization enthalpies, factors affecting IE (nuclear charge, shielding effect, penetration, electronic configuration stability), anomalies across period 2 and 3 (B vs Be, N vs O).
  \item Electron gain enthalpy: Trends across periods and down groups, anomalies (F vs Cl, O vs S).
  \item Electronegativity: Pauling scale, trends; valence and oxidation states, diagonal relationships, anomalous behaviour of second period elements.
\end{itemize}

\subsubsection{Unit 10: p-Block Elements (Groups 13 to 18)}
\begin{itemize}
  \item General introduction, electronic configuration, trends in physical and chemical properties across groups 13, 14, 15, 16, 17, and 18.
  \item Inert pair effect and its manifestation down the groups (stability of $\text{Tl}^+$ vs $\text{Tl}^{3+}$, $\text{Pb}^{2+}$ vs $\text{Pb}^{4+}$, $\text{Bi}^{3+}$ vs $\text{Bi}^{5+}$).
  \item Anomalous properties of first elements of each group (B, C, N, O, F).
  \item Allotropy: Allotropes of Carbon (diamond, graphite, fullerenes), phosphorus (white, red, black), sulfur ($\alpha, \beta$, plastic).
  \item Important trends in hydrides, oxides, and halides of p-block elements.
\end{itemize}

\subsubsection{Unit 11: d- and f-Block Elements}
\begin{itemize}
  \item \textbf{Transition Elements (3d series):} General electronic configuration, metallic character, melting and boiling points, enthalpy of atomization.
  \item Trends in 3d series: Atomic/ionic radii, ionization enthalpies, variable oxidation states, standard electrode potentials ($E^\circ_{\text{M}^{2+}/\text{M}}$ and $E^\circ_{\text{M}^{3+}/\text{M}^{2+}}$).
  \item Distinct properties: Magnetic properties (spin-only formula $\mu = \sqrt{n(n+2)}\text{ BM}$), catalytic properties, formation of coloured ions ($d-d$ transitions), interstitial compound formation, alloy formation.
  \item Compounds of Transition Metals: Preparation, properties, and oxidizing actions of Potassium Dichromate ($\text{K}_2\text{Cr}_2\text{O}_7$) and Potassium Permanganate ($\text{KMnO}_4$) in acidic, neutral, and basic media.
  \item \textbf{Lanthanoids:} Electronic configuration, oxidation states (+3 dominant, +2 and +4 states), Lanthanoid contraction and its consequences on 4d/5d radii similarity.
  \item \textbf{Actinoids:} Electronic configuration, oxidation states, comparison with lanthanoids.
\end{itemize}

\subsubsection{Unit 12: Coordination Compounds}
\begin{itemize}
  \item Werner's theory of coordination compounds (primary and secondary valencies).
  \item Coordination terminology: Coordination entity, central atom, ligands (denticity: mono-, bi-, poly-dentate, chelating ligands, ambidentate ligands), coordination number, coordination sphere, coordination polyhedron, oxidation number.
  \item IUPAC nomenclature of mononuclear coordination compounds.
  \item Isomerism: Structural isomerism (ionization, solvate/hydrate, linkage, coordination) and Stereoisomerism (geometrical cis/trans in square planar and octahedral complexes; optical isomerism in octahedral complexes).
  \item Bonding Theories:
  \begin{itemize}
    \item Valence Bond Theory (VBT): Inner and outer orbital complexes, hybridization ($sp^3, dsp^2, sp^3d^2, d^2sp^3$), magnetic properties, limitations.
    \item Crystal Field Theory (CFT): Crystal field splitting in octahedral ($\Delta_o$, $t_{2g}$ and $e_g$) and tetrahedral ($\Delta_t = \frac{4}{9}\Delta_o$) complexes, spectrochemical series, pairing energy ($P$), high spin vs low spin complexes, colour in coordination compounds.
  \end{itemize}
  \item Stability of coordination compounds and importance (biological: hemoglobin, chlorophyll, vitamin B12; analytical: qualitative analysis, water softening).
\end{itemize}

\subsection{Organic Chemistry}

\subsubsection{Unit 13: Purification and Characterisation of Organic Compounds}
\begin{itemize}
  \item Purification methods: Sublimation, crystallization, fractional distillation, steam distillation, vacuum distillation, chromatography (column, thin-layer, paper).
  \item Qualitative analysis: Lassaigne's test for detection of Nitrogen, Sulphur, and Halogens.
  \item Quantitative analysis: Estimation of carbon and hydrogen (Liebig's method), nitrogen (Dumas and Kjeldahl's methods), halogens (Carius method), sulphur, and phosphorus.
  \item Calculation of empirical and molecular formulas from percentage composition.
\end{itemize}

\subsubsection{Unit 14: Some Basic Principles of Organic Chemistry (GOC)}
\begin{itemize}
  \item IUPAC nomenclature of acyclic, cyclic, and aromatic organic compounds containing monofunctional and polyfunctional groups.
  \item Isomerism: Structural isomerism (chain, position, functional, metamerism, tautomerism: keto-enol) and Stereoisomerism (geometrical: cis/trans, E/Z; optical: chirality, enantiomers, diastereomers, meso compounds, specific rotation).
  \item Electronic Displacements: Inductive effect ($+I, -I$), Electromeric effect ($+E, -E$), Resonance/Mesomeric effect ($+M, -M$), Hyperconjugation (Baker-Nathan effect).
  \item Reactive Intermediates: Structure, stability, and generation of Carbocations, Carbanions, and Carbon Free Radicals.
  \item Reaction Types: Homolytic and heterolytic fission, electrophiles and nucleophiles, inductive vs resonance dominance in acidity/basicity of organic acids/bases.
\end{itemize}

\subsubsection{Unit 15: Hydrocarbons}
\begin{itemize}
  \item \textbf{Alkanes:} Nomenclature, isomerism, conformations of ethane (Sawhorse and Newman projections); Preparation (Wurtz reaction, Kolbe's electrolytic synthesis, decarboxylation); Chemical properties: Free radical halogenation (mechanism, selectivity of chlorination vs bromination), combustion, aromatization.
  \item \textbf{Alkenes:} Structure and geometrical isomerism; Preparation (dehydration of alcohols, dehydrohalogenation of alkyl halides - Saytzeff rule); Chemical reactions: Electrophilic addition of $X_2, HX$ (Markovnikov's rule, anti-Markovnikov peroxide effect / Kharasch effect), $H_2O, H_2SO_4$; Ozonolysis (reductive and oxidative); Oxidation with cold alkaline $\text{KMnO}_4$ (Baeyer's reagent) and acidic $\text{KMnO}_4$.
  \item \textbf{Alkynes:} Structure, acidity of terminal alkynes (reaction with $\text{NaNH}_2$, Tollens' reagent); Preparation (dehalogenation of vicinal dihalides); Reactions: Addition of $H_2, X_2, HX, H_2O$ (hydration in presence of $\text{Hg}^{2+}/\text{H}^+$ to form carbonyls), cyclic polymerization to benzene.
  \item \textbf{Aromatic Hydrocarbons:} Benzene: resonance, Huckel's $(4n+2)\pi$ electron rule for aromaticity; Electrophilic aromatic substitution (EAS): Nitration, Sulphonation, Halogenation, Friedel-Crafts Alkylation and Acylation; Directing influence of functional groups in mono-substituted benzene (ortho/para vs meta directors).
\end{itemize}

\subsubsection{Unit 16: Organic Compounds Containing Halogens (Haloalkanes \& Haloarenes)}
\begin{itemize}
  \item \textbf{Haloalkanes:} Nomenclature, nature of $\text{C}-X$ bond; Preparation from alcohols, alkenes, and hydrocarbon halogenation.
  \item Nucleophilic Substitution: Mechanisms of $S_N 1$ and $S_N 2$ reactions (stereochemistry: inversion of configuration in $S_N 2$, racemization in $S_N 1$, substrate reactivity order, solvent effects, leaving group ability).
  \item Elimination Reactions: E1 and E2 mechanisms, Saytzeff vs Hofmann elimination; Reaction with metals (formation of Grignard reagent, Wurtz reaction).
  \item \textbf{Haloarenes:} Low reactivity of aryl halides towards nucleophilic substitution (resonance, $sp^2$ carbon hybridization, instability of phenyl cation); Nucleophilic substitution under drastic conditions (Dow's process); Electrophilic substitution reactions of chlorobenzene; Fittig and Wurtz-Fittig reactions.
  \item Polyhalogen compounds: Environmental and physiological aspects of chloroform, iodoform, dichloromethane, freons, DDT.
\end{itemize}

\subsubsection{Unit 17: Organic Compounds Containing Oxygen}
\begin{itemize}
  \item \textbf{Alcohols:} Classification; Preparation (hydration of alkenes, hydroboration-oxidation, oxymercuration-demercuration, reduction of aldehydes/ketones/esters, Grignard reagents); Properties: Acidity of alcohols, esterification, reaction with $HX, PCl_3, PCl_5, SOCl_2$, dehydration to alkenes/ethers; Oxidation (PCC, Jones reagent, acidic $\text{KMnO}_4$, catalytic dehydrogenation over heated Cu at $300^\circ\text{C}$); Distinction tests (Lucas test, Victor Meyer test).
  \item \textbf{Phenols:} Preparation (from cumene, diazonium salts, chlorobenzene); Acidity of phenol vs alcohols, resonance stabilization of phenoxide ion, substituent effects; Reactions: Electrophilic aromatic substitution (bromination with $\text{Br}_2/\text{H}_2\text{O}$ vs $\text{Br}_2/\text{CS}_2$, nitration with dilute vs concentrated $\text{HNO}_3$); Kolbe's reaction, Reimer-Tiemann reaction, reaction with Zinc dust, oxidation to benzoquinone.
  \item \textbf{Ethers:} Nomenclature, Williamson's synthesis; Cleavage of $\text{C}-\text{O}$ bond by $HX$ (reaction with cold vs hot concentrated $HI$, mechanism for primary vs tertiary alkyl groups).
  \item \textbf{Aldehydes and Ketones:} Preparation (oxidation of alcohols, ozonolysis, Stephen's reaction, Rosenmund reduction, Etard reaction, Gattermann-Koch reaction, Friedel-Crafts acylation).
  \item Chemical Reactions: Nucleophilic addition reactions (addition of $HCN, NaHSO_3$, Grignard reagents, alcohols); Nucleophilic addition-elimination with ammonia derivatives ($NH_2-G$: hydroxylamine, hydrazine, 2,4-DNP, semicarbazide); Reduction: Clemmensen reduction ($\text{Zn-Hg}/\text{HCl}$), Wolff-Kishner reduction ($\text{NH}_2\text{NH}_2/\text{KOH}/\text{glycol}$), $\text{NaBH}_4, \text{LiAlH}_4$.
  \item Oxidation: Tollens' test, Fehling's test, Haloform / Iodoform reaction (compounds containing $\text{CH}_3\text{C}=\text{O}$ or $\text{CH}_3\text{CH}(\text{OH})-$ group); Reactions involving $\alpha$-hydrogen: Aldol condensation, Cross-aldol condensation, Cannizzaro reaction (compounds lacking $\alpha$-hydrogen).
  \item \textbf{Carboxylic Acids:} Acidity and substituent effects (inductive and resonance effects on $\text{p}K_a$); Preparation; Reactions: Hell-Volhard-Zelinsky (HVZ) reaction, decarboxylation (soda-lime), esterification, reduction to primary alcohols.
\end{itemize}

\subsubsection{Unit 18: Organic Compounds Containing Nitrogen}
\begin{itemize}
  \item \textbf{Amines:} Classification ($1^\circ, 2^\circ, 3^\circ$), nomenclature; Preparation (reduction of nitro compounds, nitriles, amides; Gabriel phthalimide synthesis; Hoffmann bromamide degradation reaction).
  \item Physical \& Chemical Properties: Basic character of amines (comparison in gas phase vs aqueous solution: combined inductive, steric, and solvation effects), reaction with nitrous acid ($\text{HNO}_2$), carbylamine test (isocyanide test for primary amines), Hinsberg's test (benzenesulphonyl chloride) to distinguish $1^\circ, 2^\circ$, and $3^\circ$ amines.
  \item Electrophilic aromatic substitution of aniline: Bromination ($\text{Br}_2/\text{H}_2\text{O}$ gives 2,4,6-tribromoaniline, protection by acetylation), nitration (formation of meta-nitroaniline in strongly acidic medium).
  \item \textbf{Diazonium Salts:} Preparation of benzenediazonium chloride; Chemical reactions: Sandmeyer reaction ($\text{CuCl}/\text{HCl}, \text{CuBr}/\text{HBr}, \text{CuCN}$), Gattermann reaction, Balz-Schiemann reaction ($\text{HBF}_4$ to fluorobenzene), replacement by $I, H, OH, NO_2$; Azo coupling reactions (with phenol in basic medium and aniline in slightly acidic medium to form azo dyes).
\end{itemize}

\subsubsection{Unit 19: Biomolecules}
\begin{itemize}
  \item \textbf{Carbohydrates:} Classification (monosaccharides, disaccharides, polysaccharides); Monosaccharides: Glucose (open chain structure, evidence for structure, cyclic pyranose structure, Haworth projection, $\alpha$ and $\beta$ anomers, mutarotation), Fructose (furanose structure); Disaccharides: Glycosidic linkage in Sucrose (non-reducing, invert sugar), Maltose (reducing), Lactose (reducing); Polysaccharides: Starch (amylose and amylopectin), Glycogen, Cellulose.
  \item \textbf{Proteins:} Amino acids: Classification (essential and non-essential, acidic, basic, neutral), zwitterion structure, isoelectric point; Peptide bond formation; Protein structure: Primary, Secondary ($\alpha$-helix, $\beta$-pleated sheet, hydrogen bonding), Tertiary, and Quaternary structures; Fibrous and globular proteins; Denaturation of proteins.
  \item \textbf{Enzymes:} Mechanism of enzyme action, lock and key model, induced fit, coenzymes, factors affecting enzyme activity.
  \item \textbf{Nucleic Acids:} Chemical composition: Nitrogenous bases (purines: adenine, guanine; pyrimidines: cytosine, thymine, uracil), pentose sugars (ribose, 2-deoxyribose), phosphate group; Nucleoside vs nucleotide, phosphodiester linkage; Double helical structure of DNA (Watson-Crick model, base pairing, Chargaff's rules), RNA types (mRNA, tRNA, rRNA); Biological functions (replication, protein synthesis, genetic code).
  \item \textbf{Vitamins:} Classification (fat soluble: A, D, E, K; water soluble: B-complex, C); Sources, deficiency diseases.
  \item \textbf{Hormones:} Chemical classification and elementary biological functions (steroids, peptide hormones).
\end{itemize}

\subsubsection{Unit 20: Principles Related to Practical Chemistry (NMC Specified)}
\begin{itemize}
  \item Detection of elements: Nitrogen, Sulphur, and Halogens in organic compounds.
  \item Detection of functional groups: Hydroxyl (alcoholic and phenolic), Carbonyl (aldehyde and ketone), Carboxyl, Amino groups.
  \item Inorganic Qualitative Salt Analysis:
  \begin{itemize}
    \item Cations: $\text{Pb}^{2+}, \text{Cu}^{2+}, \text{Al}^{3+}, \text{Fe}^{3+}, \text{Zn}^{2+}, \text{Ni}^{2+}, \text{Ca}^{2+}, \text{Ba}^{2+}, \text{Mg}^{2+}, \text{NH}_4^+$.
    \item Anions: $\text{CO}_3^{2-}, \text{S}^{2-}, \text{SO}_4^{2-}, \text{NO}_3^-, \text{Cl}^-, \text{Br}^-, \text{I}^-$.
  \end{itemize}
  \item Volumetric Analysis: Titration of Mohr's salt and Oxalic acid against Potassium Permanganate ($\text{KMnO}_4$), indicator and endpoint chemistry.
  \item Physical Chemistry Experiments: Enthalpy of dissolution of copper sulphate; Preparation of lyophilic and lyophobic sols.
\end{itemize}

\begin{trapbox}[Chemistry NMC Deletions \& Exclusions]
The following units and topics have been \textbf{completely deleted} by NMC:
\begin{itemize}
  \item \textbf{Physical Chemistry:} Solid State; States of Matter (Gases and Liquids); Surface Chemistry (colloids, catalysis, adsorption isotherms).
  \item \textbf{Inorganic Chemistry:} General Principles and Processes of Isolation of Elements (Metallurgy); Hydrogen; s-Block Elements (Group 1 and 2); Environmental Chemistry.
  \item \textbf{Organic Chemistry:} Polymers; Chemistry in Everyday Life.
\end{itemize}
\end{trapbox}

\newpage

% =========================================================================
% SECTION 3: BIOLOGY SYLLABUS
% =========================================================================
\section{Biology: Complete 10-Unit Detailed Breakdown}

Biology represents 50\% of the examination weightage (360 marks out of 720). The NMC syllabus covers 10 comprehensive units divided into Botany and Zoology domains.

\subsection{Unit 1: Diversity in Living World}
\begin{itemize}
  \item \textbf{The Living World:} What is living? Biodiversity, need for classification, three domains of life; Taxonomy and systematics, concept of species and taxonomical hierarchy, binomial nomenclature (ICBN, ICZN rules).
  \item \textbf{Biological Classification:} Five kingdom classification (Whittaker): Monera (Eubacteria, Archaebacteria, Cyanobacteria, Mycoplasma), Protista (Chrysophytes, Dinoflagellates, Euglenoids, Slime moulds, Protozoans), Fungi (Phycomycetes, Ascomycetes, Basidiomycetes, Deuteromycetes, lichens, mycorrhizae); Viruses, Viroids, Prions.
  \item \textbf{Plant Kingdom:} Algae (Chlorophyceae, Phaeophyceae, Rhodophyceae: pigments, food reserve, flagella, cell wall), Bryophytes (Liverworts, Mosses: gametophyte dominance, sporophyte attachment), Pteridophytes (Psilopsida, Lycopsida, Sphenopsida, Pteropsida: vascular cryptogams, heterospory, seed habit precursor), Gymnosperms (naked seeds, coralloid roots of Cycas, mycorrhizal roots of Pinus), Angiosperms (overview).
  \item \textbf{Animal Kingdom:} Salient features and classification of non-chordates up to phyla level (Porifera, Cnidaria, Ctenophora, Platyhelminthes, Aschelminthes, Annelida, Arthropoda, Mollusca, Echinodermata, Hemichordata) and chordates up to class level (Urochordata, Cephalochordata, Vertebrata: Cyclostomata, Chondrichthyes, Osteichthyes, Amphibia, Reptilia, Aves, Mammalia); Diagnostic characters: symmetry, coelom, germ layers, segmentation, notochord.
\end{itemize}

\subsection{Unit 2: Structural Organisation in Animals and Plants}
\begin{itemize}
  \item \textbf{Morphology of Flowering Plants:} Root (regions, modifications), stem (modifications), leaf (parts, venation, simple and compound leaves, phyllotaxy, modifications).
  \item \textbf{Inflorescence \& Flower:} Racemose and cymose inflorescence; Flower: parts (calyx, corolla, androecium, gynoecium), symmetry (actinomorphic, zygomorphic, asymmetric), ovary position (hypogynous, perigynous, epigynous), aestivation (valvate, twisted, imbricate, vexillary), placentation (marginal, axile, parietal, free central, basal).
  \item \textbf{Fruit \& Seed:} True fruit, false fruit, parthenocarpic fruit; Structure of dicot and monocot seed.
  \item \textbf{Detailed Plant Families (NMC Expanded Scope):}
  \begin{enumerate}
    \item \textbf{Fabaceae (Leguminosae):} Vegetative characters, floral formula, floral diagram, economic importance.
    \item \textbf{Solanaceae:} Nightshade family, floral formula, persistent calyx, epipetalous stamens.
    \item \textbf{Liliaceae:} Monocot family, perianth (tepals), trimerous flowers.
    \item \textbf{Malvaceae:} Cotton/China rose family, monadelphous stamens, epicalyx.
    \item \textbf{Cruciferae (Brassicaceae):} Mustard family, tetradynamous stamens ($2+4$), cruciform corolla.
    \item \textbf{Compositae (Asteraceae):} Sunflower family, capitulum/head inflorescence, ray and disc florets, cypsela fruit with pappus.
    \item \textbf{Gramineae (Poaceae):} Grass family, spikelet inflorescence, caryopsis fruit, feathery stigma.
  \end{enumerate}
  \item \textbf{Anatomy of Flowering Plants:} Meristematic tissues (apical, intercalary, lateral); Permanent tissues (simple: parenchyma, collenchyma, sclerenchyma; complex: xylem, phloem); Tissue systems: Epidermal, ground, and vascular tissue systems; Internal anatomy of dicotyledonous and monocotyledonous roots, stems, and leaves; Secondary growth in dicot stem and root (vascular cambium, cork cambium/phellogen, heartwood and sapwood, spring wood and autumn wood/annual rings).
  \item \textbf{Animal Tissues:} Morphology, types, and functions of Epithelial tissue (simple: squamous, cuboidal, columnar, ciliated, pseudostratified; compound: stratified), Connective tissue (loose: areolar, adipose; dense: regular [tendon, ligament], irregular; specialized: cartilage, bone, blood, lymph), Muscle tissue (skeletal, smooth, cardiac), and Nervous tissue (neurons, neuroglia).
  \item \textbf{Animal Morphology \& Anatomy (NMC Specified):}
  \begin{itemize}
    \item \textbf{Frog (\textit{Rana tigrina}):} Morphology, digestive system, respiratory system, circulatory system (three-chambered heart, portal systems), excretory system (mesonephric kidneys, ureters, cloaca), nervous system, sense organs, reproductive system (male: testes, vasa efferentia, Bidder's canal; female: ovaries, oviducts, fertilization and development).
    \item \textbf{Cockroach (\textit{Periplaneta americana}):} External morphology, digestive tract (crop, gizzard, hepatic caeca, Malpighian tubules), open circulatory system (13-chambered heart), respiratory system (spiracles, tracheae), nervous system, compound eyes (ommatidia, mosaic vision), reproductive system (male and female structures, mushroom gland, ootheca).
  \end{itemize}
\end{itemize}

\subsection{Unit 3: Cell Structure and Function}
\begin{itemize}
  \item \textbf{Cell: The Unit of Life:} Cell theory and its modern modifications (Schleiden, Schwann, Virchow: \textit{omnis cellula-e cellula}); Prokaryotic cell vs Eukaryotic cell; Cell envelope (glycocalyx, cell wall, plasma membrane), mesosomes, inclusion bodies.
  \item \textbf{Eukaryotic Cell Organelles:} Plasma membrane (Fluid Mosaic Model of Singer and Nicolson, transport mechanisms: active, passive, facilitated); Endomembrane system: Endoplasmic Reticulum (RER, SER), Golgi apparatus, Lysosomes, Vacuoles; Mitochondria (endosymbiotic origin, cristae, ATP synthesis site); Plastids (chloroplast ultrastructure, thylakoids, stroma); Ribosomes (70S vs 80S); Cytoskeleton (microfilaments, intermediate filaments, microtubules); Cilia and flagella ($9+2$ axoneme structure); Centrosome and centrioles ($9+0$ cartwheel structure); Nucleus (nuclear envelope, nuclear pores, nucleolus, chromatin structure, chromosomes, centromere, kinetochore, telomere, classification based on centromere position: metacentric, sub-metacentric, acrocentric, telocentric).
  \item \textbf{Biomolecules:} Chemical composition of living cells, biomacromolecules vs micromolecules; Carbohydrates, lipids (fatty acids, triglycerides, phospholipids - lecithin), amino acids, nucleic acids; Structure of proteins (primary, secondary, tertiary, quaternary); Dynamic state of body constituents (metabolism, concept of turnover); Enzymes: chemical nature, activation energy curve, factors affecting enzyme action (substrate concentration, Michaelis-Menten constant $K_m$, temperature, pH), enzyme inhibition (competitive, non-competitive, allosteric), classification and nomenclature of enzymes.
  \item \textbf{Cell Cycle and Cell Division:} Cell cycle phases ($G_1, S, G_2, M$ phase, $G_0$ quiescent stage); Mitosis: Prophase, Metaphase (karyotyping stage), Anaphase (chromatid separation), Telophase, Cytokinesis (cell plate vs cleavage furrow), significance of mitosis; Meiosis: Meiosis I (Prophase I substages: Leptotene, Zygotene [synapsis, synaptonemal complex], Pachytene [crossing over, recombinase enzyme], Diplotene [chiasmata, dissolution of synaptonemal complex], Diakinesis [terminalisation of chiasmata]; Metaphase I, Anaphase I [reductional segregation], Telophase I); Interkinesis; Meiosis II (equational division); Significance of meiosis in generating genetic variability.
\end{itemize}

\subsection{Unit 4: Plant Physiology}
\begin{itemize}
  \item \textbf{Photosynthesis in Higher Plants:} Photosynthesis as an autotrophic energy-converting process; Historical experiments (Priestley, Ingenhousz, Senebier, de Saussure, Engelmann, Van Niel); Photosynthetic pigments (chlorophyll $a, b$, carotenoids, xanthophylls), absorption and action spectra.
  \item \textbf{Light Reaction (Photochemical Phase):} Light harvesting complexes (LHC), Photosystem I ($P_{700}$) and Photosystem II ($P_{680}$); Photolysis of water (oxygen evolving complex, manganese and chlorine ions); Non-cyclic and cyclic photophosphorylation, Z-scheme; Chemiosmotic hypothesis (Peter Mitchell): proton gradient accumulation in thylakoid lumen, ATP synthase complex ($CF_0-CF_1$).
  \item \textbf{Dark Reaction (Biosynthetic Phase):} Calvin cycle ($C_3$ pathway): Carboxylation (RuBisCO enzyme), Reduction, Regeneration; $C_4$ pathway (Hatch-Slack pathway): Kranz anatomy, dimorphic chloroplasts, PEP carboxylase, spatial separation of carbon fixation; Photorespiration ($C_2$ cycle / glycolate pathway) in $C_3$ plants: organelles involved (chloroplast, peroxisome, mitochondrion), wastage of energy; Factors affecting photosynthesis: Blackman's law of limiting factors, light, $CO_2$ concentration, temperature, water.
  \item \textbf{Respiration in Plants:} Cellular respiration as an energy-releasing catabolic process; Glycolysis (EMP pathway): 10-step sequence, substrate-level phosphorylation, key enzymes, net gain ($2\text{ ATP} + 2\text{ NADH}$); Fermentation: Alcoholic and lactic acid fermentation.
  \item \textbf{Aerobic Respiration:} Oxidative decarboxylation of pyruvate to Acetyl-CoA (link reaction); Tricarboxylic Acid (TCA) cycle / Krebs cycle in mitochondrial matrix, steps, GTP generation; Electron Transport System (ETS) and Oxidative Phosphorylation on inner mitochondrial membrane: Complexes I, II, III, IV, and V, role of oxygen as terminal electron acceptor; Respiratory balance sheet, ATP yield; Amphibolic nature of respiratory pathway; Respiratory Quotient ($RQ$) for carbohydrates ($1.0$), fats ($0.7$), proteins ($0.9$), organic acids ($>1.0$), anaerobiosis ($ \infty $).
  \item \textbf{Plant Growth and Development:} Growth characteristics, phases of growth (meristematic, elongation, maturation), growth rates (arithmetic and geometric growth curves, sigmoid curve); Differentiation, dedifferentiation, redifferentiation; Plasticity (heterophylly in cotton, coriander, larkspur, buttercup); Plant Growth Regulators (PGRs): Discovery, physiological functions, and agricultural applications of:
  \begin{itemize}
    \item Auxins (IAA, IBA, NAA, 2,4-D: apical dominance, rooting, weedicide).
    \item Gibberellins ($GA_3$: stem elongation, bolting, malting in brewing).
    \item Cytokinins (kinetin, zeatin: cell division, delay of leaf senescence).
    \item Ethylene ($C_2H_4$: fruit ripening, triple response, femaleness).
    \item Abscisic Acid (ABA: stress hormone, stomatal closure, seed dormancy).
  \end{itemize}
\end{itemize}

\subsection{Unit 5: Human Physiology}
\begin{itemize}
  \item \textbf{Breathing and Exchange of Gases:} Respiratory organs in animals, human respiratory system anatomy; Mechanism of breathing: inspiration and expiration, role of diaphragm and intercostal muscles; Respiratory volumes and capacities (TV, IRV, ERV, RV, VC, FRC, TLC); Gas exchange at alveolar and tissue levels, partial pressures ($pO_2, pCO_2$); Transport of gases in blood: Oxygen transport (oxyhemoglobin, oxygen dissociation curve, factors shifting curve: Bohr effect), Carbon dioxide transport (as bicarbonate ions via carbonic anhydrase, carbaminohemoglobin, dissolved in plasma); Regulation of respiration (medullary rhythm centre, pneumotaxic centre, aortic and carotid chemoreceptors); Respiratory disorders: Asthma, Emphysema, Occupational Respiratory Disorders (silicosis, asbestosis).
  \item \textbf{Body Fluids and Circulation:} Blood composition: formed elements (erythrocytes, leucocytes, platelets), plasma proteins (albumin, globulin, fibrinogen); Blood groups (ABO and Rh systems, erythroblastosis fetalis); Blood coagulation cascade (thromboplastin, prothrombinase, thrombin, fibrin); Lymph (tissue fluid); Human circulatory system: internal anatomy of human heart, conducting system (SA node, AV node, bundle of His, Purkinje fibres); Cardiac cycle, heart sounds ($LUBB, DUPP$), stroke volume, cardiac output ($CO = SV \times HR$); Electrocardiogram (ECG: P wave, QRS complex, T wave); Double circulation (pulmonary and systemic, hepatic portal system); Regulation of cardiac activity (neural, autonomic, hormonal); Circulatory disorders: Hypertension, Coronary Artery Disease (Atherosclerosis), Angina Pectoris, Heart Failure.
  \item \textbf{Excretory Products and Their Elimination:} Modes of excretion (ammonotelism, ureotelism, uricotelism); Human excretory system: gross anatomy of kidney, nephron structure (cortical vs juxtamedullary nephrons), Malpighian corpuscle, PCT, loop of Henle, DCT, collecting duct; Urine formation: Glomerular filtration (GFR $\approx 125\text{ mL/min}$), tubular reabsorption, tubular secretion; Counter-current mechanism in loop of Henle and vasa recta for urine concentration; Regulation of kidney function: Renin-Angiotensin-Aldosterone System (RAAS), Antidiuretic Hormone (ADH / Vasopressin), Atrial Natriuretic Factor (ANF); Micturition reflex; Role of lungs, liver, skin in excretion; Disorders: Uremia, Renal Failure, Renal Calculi, Glomerulonephritis, Dialysis.
  \item \textbf{Locomotion and Movement:} Types of movement (amoeboid, ciliary, muscular); Skeletal muscle ultrastructure: myofibrils, sarcomere (Z-line, I-band, A-band, H-zone, M-line); Contractile proteins: Actin (F-actin, tropomyosin, troponin complex) and Myosin (heavy and light meromyosin, cross-bridges); Sliding filament mechanism of muscle contraction (role of $Ca^{2+}$, troponin binding, ATP hydrolysis); Muscle twitch, tetany, fatigue; Human skeletal system: Axial skeleton (80 bones: skull, vertebral column, ribs, sternum) and Appendicular skeleton (126 bones: pectoral and pelvic girdles, limb bones); Types of joints: Fibrous (synarthrosis), Cartilaginous (amphiarthrosis), Synovial (diarthrosis: ball and socket, hinge, pivot, gliding, saddle); Disorders: Myasthenia gravis, Tetany, Muscular dystrophy, Arthritis, Osteoporosis, Gout.
  \item \textbf{Neural Control and Coordination:} Neuron structure (dendrites, cell body, axon, myelinated vs unmyelinated fibres); Generation and conduction of nerve impulse: Resting potential ($Na^+/K^+$ pump), Action potential (depolarization, $Na^+$ influx), repolarization; Synaptic transmission: Electrical synapses and Chemical synapses (neurotransmitters, synaptic cleft, receptor binding); Central Nervous System: Forebrain (cerebrum, thalamus, hypothalamus), Midbrain (corpora quadrigemina), Hindbrain (cerebellum, pons, medulla oblongata); Reflex action and reflex arc; Sensory perception: Elementary structure and functioning of human Eye and human Ear.
  \item \textbf{Chemical Coordination and Integration:} Endocrine glands and hormone classification (peptide, steroid, iodothyronines, amino acid derivatives); Hypothalamus (releasing and inhibiting hormones); Pituitary gland (anterior: GH, PRL, TSH, ACTH, LH, FSH; posterior: oxytocin, vasopressin); Pineal gland (melatonin, circadian rhythm); Thyroid gland ($T_3, T_4$, thyrocalcitonin); Parathyroid gland (PTH, calcium homeostasis); Thymus gland (thymosins, cell-mediated immunity); Adrenal gland (adrenal cortex: mineralocorticoids, glucocorticoids; adrenal medulla: adrenaline, noradrenaline); Pancreas (islets of Langerhans: $\alpha$-cells glucagon, $\beta$-cells insulin); Gonads (testes: testosterone; ovaries: estrogen, progesterone); Hormones of heart (ANF), kidney (erythropoietin), gastrointestinal tract (gastrin, secretin, CCK, GIP); Mechanism of hormone action: Membrane-bound receptors and second messengers ($cAMP, IP_3, Ca^{2+}$); Intracellular receptors (nuclear hormone-receptor complex, gene expression regulation); Endocrine disorders: Dwarfism, Acromegaly, Diabetes Insipidus, Goitre, Cretinism, Graves' disease, Addison's disease, Cushing's syndrome, Diabetes Mellitus.
\end{itemize}

\subsection{Unit 6: Reproduction}
\begin{itemize}
  \item \textbf{Sexual Reproduction in Flowering Plants:} Flower structure; Microsporogenesis: anther wall layers (epidermis, endothecium, middle layers, tapetum), pollen sac, pollen grain structure (exine, intine, sporopollenin, germ pore, vegetative and generative cells); Megasporogenesis: ovule types, integuments, nucellus, functional megaspore, development of female gametophyte (embryo sac: 7-celled, 8-nucleate monosporic Polygonum type); Pollination: Autogamy, geitonogamy, xenogamy; Agents of pollination (anemophily, hydrophily, entomophily, ornithophily); Outbreeding devices to prevent inbreeding depression; Pollen-pistil interaction, pollen tube growth; Double fertilization: Syngamy and Triple fusion (primary endosperm nucleus, PEN); Post-fertilization events: Endosperm development (free nuclear vs cellular), Embryo development (globular, heart-shaped, torpedo, mature dicot vs monocot embryo with scutellum and coleoptile/coleorhiza); Seed and fruit formation, Pericarp; Apomixis and Polyembryony.
  \item \textbf{Human Reproduction:} Male reproductive system: Testes (scrotum, seminiferous tubules, Leydig cells, Sertoli cells), accessory ducts (rete testis, vasa efferentia, epididymis, vas deferens, ejaculatory duct, urethra), accessory glands (seminal vesicles, prostate, bulbourethral/Cowper's glands), semen composition; Female reproductive system: Ovaries, Fallopian tubes (infundibulum, ampulla, isthmus), Uterus (endometrium, myometrium, perimetrium), Cervix, Vagina, External genitalia (vulva), Mammary glands; Gametogenesis: Spermatogenesis (spermatogonia, primary spermatocyte, secondary spermatocyte, spermatid, spermatozoon; hormonal control: GnRH, LH, FSH, inhibin) and Oogenesis (oogonia, primary oocyte arrested in diplotene of prophase I, secondary oocyte arrested in metaphase II until fertilization, polar bodies); Menstrual cycle: Menstrual phase, Follicular/proliferative phase, Ovulatory surge (LH peak at day 14), Luteal/secretory phase (corpus luteum, progesterone); Fertilization, acrosome reaction, cortical reaction (block to polyspermy), cleavage, blastocyst formation, implantation; Pregnancy and embryonic development, Placenta functions (endocrine: hCG, hPL, estrogen, progesterone, relaxin); Parturition (foetal ejection reflex, oxytocin action), Lactation (colostrum, prolactin).
  \item \textbf{Reproductive Health:} Reproductive health needs and reproductive health problems; Population explosion and birth control: Natural/traditional methods (periodic abstinence, coitus interruptus, lactational amenorrhea), Barrier methods (condoms, diaphragms, cervical caps), Intrauterine Devices (IUDs: non-medicated, copper-releasing CuT/Cu7/Multiload, hormone-releasing Progestasert/LNG-20), Oral contraceptives (pills: combined, Saheli - centchroman), Injectables and implants, Surgical methods (vasectomy, tubectomy); Medical Termination of Pregnancy (MTP Act guidelines); Sexually Transmitted Infections (STIs: AIDS, Gonorrhea, Syphilis, Genital herpes, Chlamydiasis, Hepatitis-B, Trichomoniasis); Infertility and Assisted Reproductive Technologies (ART): IVF-ET (In Vitro Fertilization and Embryo Transfer: ZIFT, IUT), GIFT (Gamete Intra-Fallopian Transfer), ICSI (Intra-Cytoplasmic Sperm Injection), Artificial Insemination (AI / IUI).
\end{itemize}

\subsection{Unit 7: Genetics and Evolution}
\begin{itemize}
  \item \textbf{Mendelian Genetics \& Variation:} Mendel's experiments on garden pea (\textit{Pisum sativum}), monohybrid cross, Law of Segregation (purity of gametes), dihybrid cross, Law of Independent Assortment; Non-Mendelian inheritance: Incomplete dominance (flower colour in \textit{Mirabilis jalapa}, \textit{Antirrhinum majus}), Co-dominance (ABO blood groups in humans), Multiple alleles, Pleiotropy (phenylketonuria), Polygenic inheritance (human skin colour, kernel colour in wheat).
  \item \textbf{Chromosomal Theory \& Linkage:} Chromosomal theory of inheritance (Sutton and Boveri); Morgan's experiments on \textit{Drosophila melanogaster}: Linkage and crossing over, complete and incomplete linkage, genetic maps (Alfred Sturtevant); Sex determination mechanisms: XX-XY (human, Drosophila), ZZ-ZW (birds), XX-XO (grasshopper), Haplodiploidy (honeybee); Mutation: Gene mutations (point mutation: sickle cell anemia; frameshift mutation), chromosomal aberrations.
  \item \textbf{Genetic Disorders:} Pedigree analysis; Mendelian disorders (Autosomal dominant: Myotonic dystrophy, Huntington's chorea; Autosomal recessive: Sickle cell anemia, Cystic fibrosis, Phenylketonuria, Thalassemia; Sex-linked recessive: Hemophilia, Colour blindness); Chromosomal disorders: Aneuploidy (Down syndrome - trisomy 21, Klinefelter syndrome - 47, XXY, Turner syndrome - 45, XO), Polyploidy.
  \item \textbf{Molecular Basis of Inheritance:} DNA as genetic material: Transforming principle (Griffith), Biochemical characterisation (Avery, MacLeod, McCarty), Blender experiment (Hershey and Chase); RNA as genetic material (TMV); Structure of DNA and RNA; Packaging of DNA: Nucleosome model (histone octamer, $H_1$ histone, solenoid), euchromatin vs heterochromatin; DNA Replication: Semi-conservative nature, Meselson and Stahl's experiment (heavy nitrogen isotope $^{15}\text{N}$), Taylor's experiment on \textit{Vicia faba}, replication fork, replication machinery (DNA-dependent DNA polymerase, helicase, topoisomerase, primase, ligase, Okazaki fragments, continuous and discontinuous synthesis).
  \item \textbf{Gene Expression \& Regulation:} Transcription: Transcription unit (promoter, structural gene, terminator), template vs coding strand; Transcription in prokaryotes vs eukaryotes (RNA polymerase I, II, III; post-transcriptional modifications: capping with 7-methylguanosine, polyadenylation / tailing, splicing of introns by spliceosome); Genetic code: Salient features (triplet, degenerate, unambiguous, non-overlapping, universal, start codon AUG, stop codons UAA, UAG, UGA), Wobble hypothesis; Translation: Charging of tRNA (aminoacylation), structure of tRNA (clover leaf model), initiation, elongation, peptide bond synthesis by peptidyl transferase (23S rRNA in bacteria), termination; Regulation of gene expression: Operon concept (Jacob and Monod), Lac Operon (structural genes $z, y, a$, repressor protein, inducer lactose/allolactose, negative and positive control); Human Genome Project (HGP): Goals, methodologies (ESTs and Sequence Annotation), salient features of human genome, DNA fingerprinting (Alec Jeffreys, VNTRs, southern blot hybridization, applications in forensics and paternity testing).
  \item \textbf{Evolution:} Origin of life: Chemical evolution (Oparin-Haldane hypothesis), Urey-Miller experiment; Evidence for evolution: Paleontological evidence (fossils), comparative anatomy and morphology (homologous organs - divergent evolution; analogous organs - convergent evolution; adaptive radiation: Darwin's finches, Australian marsupials); Molecular evidence; Theories of evolution: Lamarckism, Darwinian theory (natural selection, survival of the fittest, overproduction, struggle for existence), Mutation theory of Hugo de Vries (saltation - single step large mutation); Modern synthetic theory; Hardy-Weinberg Principle: Genetic equilibrium, algebraic expression ($p^2 + 2pq + q^2 = 1$), factors disturbing equilibrium (gene flow, genetic drift - Founder effect and bottleneck effect, mutation, genetic recombination, natural selection: stabilizing, directional, disruptive); Human Evolution: Dryopithecus, Ramapithecus, Australopithecus, \textit{Homo habilis}, \textit{Homo erectus}, Neanderthal man, \textit{Homo sapiens} (cranial capacities and milestones).
\end{itemize}

\subsection{Unit 8: Biology and Human Welfare}
\begin{itemize}
  \item \textbf{Human Health and Diseases:} Concept of health and disease; Pathogens and infectious diseases in humans:
  \begin{itemize}
    \item Bacterial: Typhoid (\textit{Salmonella typhi}, Widal test), Pneumonia (\textit{Streptococcus pneumoniae}, \textit{Haemophilus influenzae}).
    \item Viral: Common cold (Rhinovirus).
    \item Protozoan: Malaria (\textit{Plasmodium vivax, P. falciparum}: lifecycle in female \textit{Anopheles} mosquito and human host, sporozoites, merozoites, hemozoin toxin, gametocytes), Amoebiasis (\textit{Entamoeba histolytica}).
    \item Helminthic: Ascariasis (\textit{Ascaris lumbricoides}), Filariasis / Elephantiasis (\textit{Wuchereria bancrofti}, vector Culex mosquito).
    \item Fungal: Ringworm (\textit{Microsporum, Trichophyton, Epidermophyton}).
  \end{itemize}
  \item \textbf{Immunity:} Innate immunity (physical, physiological, cellular, cytokine barriers); Acquired immunity (active and passive, natural and artificial); Humoral immune response (B-lymphocytes, antibody structure: $H_2 L_2$, classes: IgG, IgA, IgM, IgE, IgD) and Cell-mediated immune response (T-lymphocytes, graft rejection); Vaccination and immunization; Allergies (histamine, serotonin from mast cells, IgE); Autoimmunity (Rheumatoid arthritis, Hashimoto's disease); Immune system organs: Primary lymphoid organs (bone marrow, thymus) and Secondary lymphoid organs (spleen, lymph nodes, tonsils, Peyer's patches, MALT).
  \item \textbf{Non-Infectious Diseases:} Cancer: Benign vs malignant tumours, metastasis, oncogenic viruses, proto-oncogenes, carcinogens, diagnosis and treatment (surgery, radiotherapy, chemotherapy, immunotherapy); AIDS: Human Immunodeficiency Virus (HIV, retrovirus structure, reverse transcriptase), transmission, pathogenesis (attack on macrophages and helper T-cells $T_H$), ELISA test, prevention and control (NACO).
  \item \textbf{Drugs and Alcohol Abuse:} Opioids (morphine, heroin/smack, receptors in CNS and GI tract), Cannabinoids (marijuana, hashish, charas, ganja, cardiovascular effects), Coca alkaloids/Cocaine (coca plant, dopamine interference, euphoria, hallucinations), Hallucinogens (LSD, \textit{Atropa belladonna}, \textit{Datura}); Addiction, dependence, withdrawal syndrome, adolescence and abuse prevention.
  \item \textbf{Microbes in Human Welfare:} In household food processing: Curd formation (\textit{Lactobacillus} / LAB, vitamin B12 enrichment), dough fermentation (Baker's yeast \textit{Saccharomyces cerevisiae}), cheese production (Swiss cheese \textit{Propionibacterium shermanii}, Roquefort cheese); Industrial products: Fermented beverages (distilled vs undistilled), Antibiotics (Penicillin - Alexander Fleming, Florey and Chain); Chemicals, enzymes, and bioactive molecules: Citric acid (\textit{Aspergillus niger}), Acetic acid (\textit{Acetobacter aceti}), Butyric acid (\textit{Clostridium butylicum}), Lactic acid (\textit{Lactobacillus}), Lipases, Pectinases, Proteases, Streptokinase (clot buster from \textit{Streptococcus}), Cyclosporin-A (immunosuppressive agent from \textit{Trichoderma polysporum}), Statins (blood cholesterol lowering agent from \textit{Monascus purpureus}); Sewage treatment: Primary treatment (physical filtration and sedimentation) and Secondary biological treatment (aerobic aeration tank, flocs, BOD reduction, anaerobic sludge digester, biogas production); Biogas production: Methanogens (\textit{Methanobacterium}); Microbes as Biocontrol Agents: \textit{Bacillus thuringiensis} (Bt), Ladybird beetle, Dragonflies, \textit{Trichoderma}, Baculoviruses (Nucleopolyhedrovirus); Microbes as Biofertilizers: Rhizobium, Azotobacter, Azospirillum, Mycorrhiza (Glomus), Cyanobacteria (Anabaena, Nostoc, Oscillatoria).
\end{itemize}

\subsection{Unit 9: Biotechnology and Its Applications}
\begin{itemize}
  \item \textbf{Principles and Processes of Biotechnology:} Definition of biotechnology (EFB definition); Genetic engineering and chemical engineering; Recombinant DNA (rDNA) technology tools:
  \begin{itemize}
    \item Restriction Enzymes: Endonucleases vs exonucleases, palindromic recognition sequences, sticky ends vs blunt ends, nomenclature (e.g. EcoRI).
    \item Cloning Vectors: Characteristics of good vector (origin of replication \textit{ori}, selectable markers: $amp^R, tet^R$, insertional inactivation / blue-white selection using $\beta$-galactosidase gene \textit{lacZ}, cloning sites); Types of vectors (pBR322, Ti plasmid of \textit{Agrobacterium tumefaciens}, retroviruses in animals).
    \item Competent Host: Methods of DNA uptake (divalent cations $\text{Ca}^{2+}$, heat shock, microinjection, gene gun / biolistics, disarmed pathogens).
  \end{itemize}
  \item \textbf{Processes of Recombinant DNA Technology:} Isolation of genetic material (DNA extraction using lysozyme, cellulase, chitinase, RNAse, protease; chilled ethanol precipitation), cutting of DNA at specific locations, agarose gel electrophoresis (separation based on size, ethidium bromide staining, UV visualization, elution); Amplification of gene of interest by Polymerase Chain Reaction (PCR: Denaturation at $\sim 94^\circ\text{C}$, Annealing at $\sim 54^\circ\text{C}$, Extension at $\sim 72^\circ\text{C}$ using Taq polymerase from \textit{Thermus aquaticus}); Insertion of rDNA into host; Culturing host cells: Bioreactors (stirred-tank bioreactors, sparged stirred-tank bioreactor, aeration, pH and temp control); Downstream processing (separation, purification, clinical trials, quality testing).
  \item \textbf{Biotechnology and Its Applications:}
  \begin{itemize}
    \item Applications in Agriculture: Genetically Modified Organisms (GMOs); Bt crops: Bt cotton (Cry toxins: \textit{cryIAc}, \textit{cryIIAb} against cotton bollworms, \textit{cryIAb} against corn borer; activation of protoxin in alkaline insect midgut); Pest-resistant plants: RNA interference (RNAi) in tobacco against nematode \textit{Meloidogyne incognita} (cellular defense using double-stranded RNA, gene silencing).
    \item Applications in Medicine: Genetically engineered insulin (Humulin: Eli Lilly, synthesis of chains A and B separately linked by disulfide bonds, elimination of C-peptide); Gene therapy: First clinical gene therapy (1990) for ADA (adenosine deaminase) deficiency in 4-year-old girl; Molecular diagnostics: PCR, ELISA, recombinant DNA probes for early pathogen and genetic defect detection.
    \item Transgenic Animals: Transgenic mice, sheep, cows (Rosie cow: human protein-enriched milk containing $\alpha$-lactalbumin); Applications in testing drug safety, vaccine safety (polio vaccine in transgenic mice), chemical testing, and disease models.
    \item Ethical Issues: GEAC (Genetic Engineering Appraisal Committee); Biopiracy (exploitation of indigenous bioresources without authorization: Indian Basmati rice, Neem, Turmeric patent disputes).
  \end{itemize}
\end{itemize}

\subsection{Unit 10: Ecology and Environment}
\begin{itemize}
  \item \textbf{Organisms and Populations:} Organism and its environment: abiotic factors (temperature, water, light, soil); Responses to abiotic factors: Regulators, conformers, partial regulators, migration, suspension (spores, dormancy, hibernation in bears, aestivation in snails, diapause in zooplankton); Adaptations: Morphological, physiological, and behavioural (Kangaroo rat water conservation, desert plants CAM pathway and sunken stomata, Allen's rule in cold climate mammals, high altitude sickness acclimation).
  \item \textbf{Population Ecology:} Population attributes: Natality, mortality, immigration, emigration, sex ratio, age pyramids (expanding, stable, declining); Population growth models: Exponential growth ($\frac{dN}{dt} = rN$, $N_t = N_0 e^{rt}$) and Logistic growth ($\frac{dN}{dt} = rN[\frac{K - N}{K}]$, sigmoid curve, carrying capacity $K$); Population interactions:
  \begin{itemize}
    \item Mutualism (+/+): Lichens, mycorrhizae, fig and wasp, pollination syndromic adaptations.
    \item Competition (-/-): Gause's competitive exclusion principle, resource partitioning (MacArthur's warbler birds).
    \item Predation (+/-): Biological control, maintenance of species diversity (Pisaaster starfish experiment).
    \item Parasitism (+/-): Ectoparasites, endoparasites, brood parasitism in cuckoo/koel.
    \item Commensalism (+/0): Orchid on mango branch, barnacles on whale, cattle egret and grazing cattle, clown fish and sea anemone.
    \item Amensalism (-/0): Penicillium and bacteria.
  \end{itemize}
  \item \textbf{Ecosystem:} Structure and function: Productivity (Primary productivity: GPP, NPP = $\text{GPP} - R$; Secondary productivity); Decomposition: Fragmentation, leaching, catabolism, humification, mineralization; Energy flow: Unidirectional flow, 10\% law of Lindeman, food chain (grazing vs detritus food chain), food web; Ecological pyramids: Pyramid of numbers (upright and inverted in tree ecosystem), pyramid of biomass (upright on land, inverted in aquatic ecosystems), pyramid of energy (always upright, second law of thermodynamics).
  \item \textbf{Biodiversity and Its Conservation:} Concept of biodiversity, levels: Genetic diversity (\textit{Rauwolfia vomitoria} reserpine, rice varieties, mango strains), Species diversity (Western Ghats amphibians), Ecological diversity (India's biomes); Global species estimates (Robert May's estimate: 7 million species); Patterns of biodiversity: Latitudinal gradients (species richness highest in tropics, Amazon rainforest diversity), Species-Area relationship (Alexander von Humboldt: $\log S = \log C + Z\log A$, values of $Z$); Importance of biodiversity: David Tilman's field experiments, Rivet Popper hypothesis (Paul Ehrlich); Loss of biodiversity: Evil Quartet (Habitat loss and fragmentation, Over-exploitation, Alien species invasions [Nile perch, water hyacinth, Parthenium, African catfish], Co-extinctions); Conservation strategies: In-situ conservation (National Parks, Wildlife Sanctuaries, Biosphere Reserves, Sacred Groves, Biodiversity Hotspots: Western Ghats-Sri Lanka, Indo-Burma, Himalaya) and Ex-situ conservation (Zoological parks, Botanical gardens, Wildlife safari parks, Cryopreservation, Seed banks, Tissue culture).
\end{itemize}

\begin{trapbox}[Biology NMC Deletions \& Exclusions]
The following chapters have been \textbf{completely deleted} by NMC from the NEET examination:
\begin{itemize}
  \item \textbf{Plant Physiology:} Transport in Plants (xylem/phloem translocation, water potential, plasmolysis); Mineral Nutrition (essential macro/micro-elements, deficiency symptoms, nitrogen cycle).
  \item \textbf{Human Physiology:} Digestion and Absorption (digestive tract anatomy, enzyme action, absorption of carbohydrates/fats/proteins, digestive disorders).
  \item \textbf{Botany/Zoology General:} Strategies for Enhancement in Food Production (plant breeding, tissue culture, animal husbandry); Environmental Issues (air, water, solid waste pollution, greenhouse effect, ozone depletion, deforestation).
\end{itemize}
\end{trapbox}

\newpage

% =========================================================================
% SECTION 4: SYLLABUS MAPPING \& HIGH-YIELD WEIGHTAGE
% =========================================================================
\section{Syllabus Mapping \& Empirical NEET Question Distribution}

To maximize study efficiency, this section categorizes all 50 units by historical question frequency and typical NEET question weightage.

\subsection{Subject-Wise Weightage Summary}

\begin{table}[htbp]
\centering
\small
\renewcommand{\arraystretch}{1.2}
\begin{tabularx}{\textwidth}{L c c Y}
\toprule
\textbf{Subject Domain} & \textbf{Total Units} & \textbf{Typical Questions} & \textbf{Marks Allotted} \\
\midrule
\textbf{Physics: Mechanics \& Thermal} & Units 1--10 & 22--24 Questions & 88--96 Marks \\
\textbf{Physics: Electrodynamics \& Modern} & Units 11--20 & 21--23 Questions & 84--92 Marks \\
\midrule
\textbf{Chemistry: Physical Chemistry} & Units 1--8 & 15--17 Questions & 60--68 Marks \\
\textbf{Chemistry: Inorganic Chemistry} & Units 9--12 & 13--15 Questions & 52--60 Marks \\
\textbf{Chemistry: Organic \& Practical} & Units 13--20 & 14--16 Questions & 56--64 Marks \\
\midrule
\textbf{Botany: Genetics, Diversity, Ecology} & 5 Domains & 45 Questions & 180 Marks \\
\textbf{Zoology: Physiology, Reproduction, Evolution}& 5 Domains & 45 Questions & 180 Marks \\
\midrule
\textbf{Master NEET Total} & \textbf{50 Units} & \textbf{180 Questions (All Compulsory)} & \textbf{720 Marks} \\
\bottomrule
\end{tabularx}
\caption{Empirical Question Breakdown Across Official NEET UG Units.}
\end{table}

\subsection{Tier 1 High-Yield Target Units (Mandatory Top Priority)}

\begin{enumerate}
  \item \textbf{Physics:}
  \begin{itemize}
    \item Current Electricity \& Circuits (3--4 Questions)
    \item Ray \& Wave Optics (4--5 Questions)
    \item Electrostatics \& Capacitors (3--4 Questions)
    \item Modern Physics (Dual Nature, Atoms, Nuclei) (4--5 Questions)
    \item Laws of Motion, Work-Energy \& Rotational Dynamics (5--6 Questions)
  \end{itemize}
  \item \textbf{Chemistry:}
  \begin{itemize}
    \item Chemical Bonding \& Molecular Structure (3--4 Questions)
    \item Coordination Compounds (3--4 Questions)
    \item Equilibrium (Chemical \& Ionic) (3--4 Questions)
    \item Organic Reaction Mechanisms (GOC, Carbonyls, Amines) (7--9 Questions)
    \item Electrochemistry \& Chemical Kinetics (4--5 Questions)
  \end{itemize}
  \item \textbf{Biology:}
  \begin{itemize}
    \item Molecular Basis of Inheritance \& Mendelian Genetics (14--18 Questions)
    \item Human Physiology: Circulatory, Excretory, Neural, Endocrine (12--15 Questions)
    \item Plant Physiology: Photosynthesis \& Respiration (6--8 Questions)
    \item Cell Biology \& Cell Division (7--9 Questions)
    \item Biotechnology: Principles \& Applications (7--9 Questions)
    \item Ecology: Populations, Ecosystems, Biodiversity (8--10 Questions)
  \end{itemize}
\end{enumerate}

\vspace{15pt}
\begin{center}
  \rule{0.6\textwidth}{0.8pt}\\[6pt]
  {\small\textbf{\color{primarydark}End of NEET UG 2027 Master Syllabus Specification}}\\[2pt]
  {\footnotesize\color{accentblue}Prepared by Satya Prakash \quad\textbar\quad Strict NCERT \& NMC Benchmark}
\end{center}

\end{document}
